\documentclass{article}
\usepackage{amsmath,amssymb,amsthm}
\usepackage{algorithm,algorithmic}
\usepackage{bm}
\usepackage{hyperref}
\newtheorem{theorem}{Theorem}
\newtheorem{lemma}{Lemma}
\newtheorem{assumption}{Assumption}
\newtheorem{definition}{Definition}
\newtheorem{remark}{Remark}
\begin{document}

\section*{Differentially Private Stochastic Gradient Descent (DP‑SGD) with Gaussian Noise}

\section*{Mathematical Formulation}
Let $f:\mathbb{R}^d\rightarrow\mathbb{R}$ be the (possibly non‑convex) loss function defined as the expectation over a data distribution $\mathcal{D}$,
\[
f(\mathbf{w})=\mathbb{E}_{\xi\sim\mathcal{D}}\big[\ell(\mathbf{w};\xi)\big].
\]
At iteration $t\in\{0,1,\dots,T-1\}$ we sample a minibatch $\mathcal{B}_t$ of size $b$,
\[
\mathbf{g}_t = \frac{1}{b}\sum_{\xi\in\mathcal{B}_t}\operatorname{clip}_C\big(\nabla\ell(\mathbf{w}_t;\xi)\big) \;+\; \bm{\zeta}_t,
\qquad
\bm{\zeta}_t\sim\mathcal{N}(\mathbf{0},\sigma^2\mathbf{I}_d).
\]
The clipping operator is defined component‑wise as
\[
\operatorname{clip}_C(\mathbf{u})=
\begin{cases}
\displaystyle \frac{C}{\|\mathbf{u}\|_2}\,\mathbf{u}, & \text{if }\|\mathbf{u}\|_2>C,\\[4pt]
\mathbf{u}, & \text{otherwise.}
\end{cases}
\]
The parameter update is
\[
\boxed{\;\mathbf{w}_{t+1}= \mathbf{w}_t-\eta_t\,\mathbf{g}_t\;}
\tag{1}
\]
where $\eta_t>0$ is the learning‑rate schedule. The Gaussian noise $\bm{\zeta}_t$ ensures $(\varepsilon,\delta)$‑differential privacy; the privacy accountant (e.g. moments accountant) tracks the cumulative privacy loss over $T$ iterations.

\section*{Pseudocode}
\begin{algorithm}[H]
\caption{DP‑SGD with Gaussian Noise}
\begin{algorithmic}[1]
\STATE \textbf{Input:} initial point $\mathbf{w}_0$, learning rates $\{\eta_t\}_{t=0}^{T-1}$, clipping bound $C$, noise scale $\sigma$, minibatch size $b$, total iterations $T$.
\FOR{$t=0$ \TO $T-1$}
    \STATE Sample a minibatch $\mathcal{B}_t$ of size $b$ uniformly from the dataset.
    \STATE Compute per‑example gradients $\mathbf{u}_\xi=\nabla\ell(\mathbf{w}_t;\xi)$ for $\xi\in\mathcal{B}_t$.
    \STATE Clip: $\tilde{\mathbf{u}}_\xi=\operatorname{clip}_C(\mathbf{u}_\xi)$.
    \STATE Form the averaged clipped gradient 
    \[
    \bar{\mathbf{g}}_t=\frac{1}{b}\sum_{\xi\in\mathcal{B}_t}\tilde{\mathbf{u}}_\xi .
    \]
    \STATE Sample Gaussian noise $\bm{\zeta}_t\sim\mathcal{N}(\mathbf{0},\sigma^2\mathbf{I}_d)$.
    \STATE Noisy gradient $\mathbf{g}_t=\bar{\mathbf{g}}_t+\bm{\zeta}_t$.
    \STATE Update $\mathbf{w}_{t+1}= \mathbf{w}_t-\eta_t\mathbf{g}_t$.
\ENDFOR
\STATE \textbf{Output:} $\mathbf{w}_T$ (or the average $\frac{1}{T}\sum_{t=0}^{T-1}\mathbf{w}_t$).
\end{algorithmic}
\end{algorithm}

\section*{Dry Run Example}
Consider the scalar quadratic $f(w)=(w-3)^2$, whose gradient is $\nabla f(w)=2(w-3)$.  
Parameters: $w_0=0$, learning rate $\eta=0.1$, clipping bound $C=10$ (no effect), noise variance $\sigma^2=0$ (to illustrate the deterministic part), and a batch size $b=1$ (single data point).  

\begin{center}
\begin{tabular}{c|c|c|c}
Iteration $t$ & $w_t$ & $g_t=\nabla f(w_t)$ & $w_{t+1}=w_t-\eta g_t$\\\hline
0 & $0$ & $2(0-3)=-6$ & $0-0.1(-6)=0.6$\\
1 & $0.6$ & $2(0.6-3)=-4.8$ & $0.6-0.1(-4.8)=1.08$\\
2 & $1.08$ & $2(1.08-3)=-3.84$ & $1.08-0.1(-3.84)=1.464$\\
\end{tabular}
\end{center}
Objective values:
\[
f(0)=9,\qquad f(0.6)=5.76,\qquad f(1.08)=3.6864,\qquad f(1.464)=2.2237.
\]
The iterates move toward the minimiser $w^\star=3$, illustrating the update (1).

\newpage
\section*{Assumptions}
\begin{assumption}[Smoothness]\label{ass:smooth}
$f$ is $L$‑smooth, i.e. for all $\mathbf{u},\mathbf{v}\in\mathbb{R}^d$,
\[
f(\mathbf{u})\le f(\mathbf{v})+\langle\nabla f(\mathbf{v}),\mathbf{u}-\mathbf{v}\rangle+\frac{L}{2}\|\mathbf{u}-\mathbf{v}\|_2^{2}.
\]
\end{assumption}
\begin{assumption}[Bounded Gradient Norm after Clipping]\label{ass:clip}
For every data point $\xi$ and any $\mathbf{w}$,
\[
\big\|\operatorname{clip}_C\big(\nabla\ell(\mathbf{w};\xi)\big)\big\|_2\le C.
\]
Consequently $\|\bar{\mathbf{g}}_t\|_2\le C$.
\end{assumption}
\begin{assumption}[Unbiased Stochastic Gradient]\label{ass:unbiased}
Conditioned on $\mathbf{w}_t$, the averaged clipped gradient satisfies
\[
\mathbb{E}\big[\bar{\mathbf{g}}_t\mid\mathbf{w}_t\big]=\nabla f(\mathbf{w}_t).
\]
\end{assumption}
\begin{assumption}[Finite Variance of Stochastic Gradient]\label{ass:variance}
There exists $G^2\ge0$ such that for all $t$,
\[
\mathbb{E}\big[\|\bar{\mathbf{g}}_t-\nabla f(\mathbf{w}_t)\|_2^{2}\mid\mathbf{w}_t\big]\le \frac{G^{2}}{b}.
\]
\end{assumption}
\begin{assumption}[Gaussian Noise]\label{ass:gauss}
The added noise $\bm{\zeta}_t\sim\mathcal{N}(\mathbf{0},\sigma^{2}\mathbf{I}_d)$ is independent of all past randomness and of the data.
\end{assumption}

\section*{Main Convergence Theorem}
\begin{theorem}[Expected Gradient Norm Bound for DP‑SGD]\label{thm:main}
Let $\{\mathbf{w}_t\}_{t=0}^{T}$ be generated by the DP‑SGD update (1) under Assumptions \ref{ass:smooth}–\ref{ass:gauss}. Choose a constant learning rate $\eta_t\equiv\eta\le\frac{1}{L}$. Then
\[
\frac{1}{T}\sum_{t=0}^{T-1}\mathbb{E}\big[\|\nabla f(\mathbf{w}_t)\|_2^{2}\big]
\;\le\;
\underbrace{\frac{2\big(f(\mathbf{w}_0)-f^\star\big)}{\eta T}}_{\text{optimization term}}
\;+\;
\underbrace{L\eta\Big(\frac{G^{2}}{b}+d\sigma^{2}\Big)}_{\text{stochastic + privacy noise term}} .
\tag{2}
\]
Consequently, by setting $\eta=\Theta\!\big(\frac{1}{\sqrt{T}}\big)$ we obtain
\[
\frac{1}{T}\sum_{t=0}^{T-1}\mathbb{E}\big[\|\nabla f(\mathbf{w}_t)\|_2^{2}\big]=O\!\Big(\frac{1}{\sqrt{T}}\Big)
\;+\;O\!\big(Ld\sigma^{2}\eta\big).
\]
\end{theorem}

\section*{Theoretical Properties}
\begin{itemize}
\item \textbf{Stability.} The clipping bound $C$ guarantees that the influence of any single data point on $\mathbf{g}_t$ is bounded, which is essential for differential privacy and also prevents exploding updates.
\item \textbf{Hyper‑parameter Sensitivity.} The bound (2) shows a linear dependence on the learning rate $\eta$ for the noise term and an inverse dependence for the optimization term; the optimal $\eta$ balances these two effects, yielding the $\mathcal{O}(1/\sqrt{T})$ rate.
\item \textbf{Comparison to Non‑private SGD.} When $\sigma=0$ (no privacy noise) and $b$ is large so that $G^{2}/b$ is negligible, (2) reduces to the classical non‑convex SGD bound $\mathcal{O}(1/\sqrt{T})$.
\item \textbf{Privacy–Utility Trade‑off.} The term $d\sigma^{2}$ captures the degradation caused by the Gaussian mechanism; tighter privacy (smaller $\varepsilon$) requires larger $\sigma$, thus inflating the bound.
\end{itemize}

\section*{Preliminaries (Definitions and Lemmas)}
\begin{definition}[Gradient Noise Decomposition]
Define the total noise $\mathbf{e}_t$ as
\[
\mathbf{e}_t\;:=\;\big(\bar{\mathbf{g}}_t-\nabla f(\mathbf{w}_t)\big)+\bm{\zeta}_t .
\]
Thus $\mathbf{g}_t=\nabla f(\mathbf{w}_t)+\mathbf{e}_t$.
\end{definition}
\begin{lemma}[Smoothness Inequality]\label{lem:smooth}
If $f$ is $L$‑smooth then for any $\mathbf{w},\mathbf{u}$
\[
f(\mathbf{u})\le f(\mathbf{w})+\langle\nabla f(\mathbf{w}),\mathbf{u}-\mathbf{w}\rangle+\frac{L}{2}\|\mathbf{u}-\mathbf{w}\|_2^{2}.
\]
\end{lemma}
\begin{proof}
Standard; follows from integrating the Lipschitz condition on $\nabla f$.
\end{proof}
\begin{lemma}[Noise Moment Bound]\label{lem:noise}
Under Assumptions \ref{ass:variance} and \ref{ass:gauss},
\[
\mathbb{E}\big[\|\mathbf{e}_t\|_2^{2}\mid\mathbf{w}_t\big]
\;=\;
\mathbb{E}\big[\|\bar{\mathbf{g}}_t-\nabla f(\mathbf{w}_t)\|_2^{2}\mid\mathbf{w}_t\big]
\;+\;
\mathbb{E}\big[\|\bm{\zeta}_t\|_2^{2}\big]
\;\le\;
\frac{G^{2}}{b}+d\sigma^{2}.
\]
\end{lemma}
\begin{proof}
Independence of $\bm{\zeta}_t$ yields $\mathbb{E}[\langle\bar{\mathbf{g}}_t-\nabla f(\mathbf{w}_t),\bm{\zeta}_t\rangle]=0$ and $\mathbb{E}\|\bm{\zeta}_t\|_2^{2}=d\sigma^{2}$ because each coordinate has variance $\sigma^{2}$.
\end{proof}

\section*{Main Proof (Step‑by‑Step Derivation)}
We prove Theorem \ref{thm:main}. All expectations are taken over the randomness of minibatch sampling and the Gaussian noise unless otherwise noted.

\noindent\textbf{Step 1 (Smoothness expansion).}  
From Lemma \ref{lem:smooth} with $\mathbf{u}=\mathbf{w}_{t+1}$ and $\mathbf{w}=\mathbf{w}_t$,
\begin{align}
f(\mathbf{w}_{t+1})
&\le f(\mathbf{w}_t)+\langle\nabla f(\mathbf{w}_t),\mathbf{w}_{t+1}-\mathbf{w}_t\rangle
      +\frac{L}{2}\|\mathbf{w}_{t+1}-\mathbf{w}_t\|_2^{2}. \tag{3}
\end{align}

\noindent\textbf{Step 2 (Plug the update).}  
Using $\mathbf{w}_{t+1}-\mathbf{w}_t=-\eta\mathbf{g}_t$,
\begin{align}
f(\mathbf{w}_{t+1})
&\le f(\mathbf{w}_t)-\eta\langle\nabla f(\mathbf{w}_t),\mathbf{g}_t\rangle
      +\frac{L\eta^{2}}{2}\|\mathbf{g}_t\|_2^{2}. \tag{4}
\end{align}

\noindent\textbf{Step 3 (Decompose $\mathbf{g}_t$).}  
Recall $\mathbf{g}_t=\nabla f(\mathbf{w}_t)+\mathbf{e}_t$. Substituting,
\begin{align}
f(\mathbf{w}_{t+1})
&\le f(\mathbf{w}_t)-\eta\|\nabla f(\mathbf{w}_t)\|_2^{2}
      -\eta\langle\nabla f(\mathbf{w}_t),\mathbf{e}_t\rangle
      +\frac{L\eta^{2}}{2}\big\|\nabla f(\mathbf{w}_t)+\mathbf{e}_t\big\|_2^{2}. \tag{5}
\end{align}

\noindent\textbf{Step 4 (Expand the squared norm).}  
\[
\big\|\nabla f(\mathbf{w}_t)+\mathbf{e}_t\big\|_2^{2}
= \|\nabla f(\mathbf{w}_t)\|_2^{2}
  +2\langle\nabla f(\mathbf{w}_t),\mathbf{e}_t\rangle
  +\|\mathbf{e}_t\|_2^{2}.
\]
Insert this into (5):
\begin{align}
f(\mathbf{w}_{t+1})
&\le f(\mathbf{w}_t)-\eta\|\nabla f(\mathbf{w}_t)\|_2^{2}
      -\eta\langle\nabla f(\mathbf{w}_t),\mathbf{e}_t\rangle \nonumber\\
&\qquad+\frac{L\eta^{2}}{2}\|\nabla f(\mathbf{w}_t)\|_2^{2}
      +L\eta^{2}\langle\nabla f(\mathbf{w}_t),\mathbf{e}_t\rangle
      +\frac{L\eta^{2}}{2}\|\mathbf{e}_t\|_2^{2}. \tag{6}
\end{align}

\noindent\textbf{Step 5 (Collect like terms).}  
Group the terms containing $\langle\nabla f(\mathbf{w}_t),\mathbf{e}_t\rangle$:
\begin{align}
f(\mathbf{w}_{t+1})
&\le f(\mathbf{w}_t)
   -\Big(\eta-\frac{L\eta^{2}}{2}\Big)\|\nabla f(\mathbf{w}_t)\|_2^{2}
   +\big(L\eta^{2}-\eta\big)\langle\nabla f(\mathbf{w}_t),\mathbf{e}_t\rangle
   +\frac{L\eta^{2}}{2}\|\mathbf{e}_t\|_2^{2}. \tag{7}
\end{align}

\noindent\textbf{Step 6 (Choose $\eta\le\frac{1}{L}$).}  
If $\eta\le\frac{1}{L}$ then $\eta-\frac{L\eta^{2}}{2}\ge\frac{\eta}{2}$. Moreover,
$L\eta^{2}-\eta\le0$, so the cross term is non‑positive in expectation (see Step 8). Hence,
\begin{align}
f(\mathbf{w}_{t+1})
&\le f(\mathbf{w}_t)
   -\frac{\eta}{2}\|\nabla f(\mathbf{w}_t)\|_2^{2}
   +\frac{L\eta^{2}}{2}\|\mathbf{e}_t\|_2^{2}. \tag{8}
\end{align}

\noindent\textbf{Step 7 (Take conditional expectation).}  
Condition on $\mathbf{w}_t$ and use Lemma \ref{lem:noise}:
\begin{align}
\mathbb{E}\big[f(\mathbf{w}_{t+1})\mid\mathbf{w}_t\big]
&\le f(\mathbf{w}_t)
   -\frac{\eta}{2}\|\nabla f(\mathbf{w}_t)\|_2^{2}
   +\frac{L\eta^{2}}{2}\,\mathbb{E}\big[\|\mathbf{e}_t\|_2^{2}\mid\mathbf{w}_t\big] \nonumber\\
&\le f(\mathbf{w}_t)
   -\frac{\eta}{2}\|\nabla f(\mathbf{w}_t)\|_2^{2}
   +\frac{L\eta^{2}}{2}\Big(\frac{G^{2}}{b}+d\sigma^{2}\Big). \tag{9}
\end{align}

\noindent\textbf{Step 8 (Eliminate the cross term).}  
We used $\mathbb{E}[\langle\nabla f(\mathbf{w}_t),\mathbf{e}_t\rangle\mid\mathbf{w}_t]=0$ because
\[
\mathbb{E}[\mathbf{e}_t\mid\mathbf{w}_t]
= \mathbb{E}[\bar{\mathbf{g}}_t-\nabla f(\mathbf{w}_t)\mid\mathbf{w}_t]
  +\mathbb{E}[\bm{\zeta}_t]=\mathbf{0},
\]
by Assumption \ref{ass:unbiased} and the zero‑mean Gaussian noise.

\noindent\textbf{Step 9 (Unconditional expectation and telescoping).}  
Take total expectation and sum (9) over $t=0,\dots,T-1$:
\begin{align}
\mathbb{E}[f(\mathbf{w}_T)]
&\le f(\mathbf{w}_0)
   -\frac{\eta}{2}\sum_{t=0}^{T-1}\mathbb{E}\big[\|\nabla f(\mathbf{w}_t)\|_2^{2}\big]
   +\frac{L\eta^{2}T}{2}\Big(\frac{G^{2}}{b}+d\sigma^{2}\Big). \tag{10}
\end{align}
Since $f(\mathbf{w}_T)\ge f^\star$, rearrange (10):
\begin{align}
\frac{1}{T}\sum_{t=0}^{T-1}\mathbb{E}\big[\|\nabla f(\mathbf{w}_t)\|_2^{2}\big]
&\le
\frac{2\big(f(\mathbf{w}_0)-f^\star\big)}{\eta T}
   +L\eta\Big(\frac{G^{2}}{b}+d\sigma^{2}\Big). \tag{11}
\end{align}
Equation (11) is exactly the bound (2) claimed in Theorem \ref{thm:main}.  ∎

\section*{Rate Derivation (With Constants)}
Define the initial optimality gap $\Delta_0:=f(\mathbf{w}_0)-f^\star$.  
Choosing a constant learning rate
\[
\eta = \min\Big\{\frac{1}{L},\sqrt{\frac{2\Delta_0}{L\big(\frac{G^{2}}{b}+d\sigma^{2}\big)T}}\Big\},
\]
the two terms in (11) become equal, yielding
\[
\frac{1}{T}\sum_{t=0}^{T-1}\mathbb{E}\big[\|\nabla f(\mathbf{w}_t)\|_2^{2}\big]
\;\le\;
2\sqrt{\frac{2L\Delta_0\big(\frac{G^{2}}{b}+d\sigma^{2}\big)}{T}}.
\]
Thus the average squared gradient norm decays as $\mathcal{O}\!\big(T^{-1/2}\big)$, with a multiplicative factor that grows linearly with the privacy‑induced noise level $\sigma\sqrt{d}$.

\section*{Special Cases}
\begin{itemize}
\item \textbf{No privacy noise ($\sigma=0$).} The bound reduces to the classical non‑convex SGD result $\mathcal{O}\!\big(\frac{1}{\sqrt{T}}\big)$ with variance term $\frac{G^{2}}{b}$ only.
\item \textbf{Full‑batch ($b=n$).} When the whole dataset is used, $G^{2}/b\approx0$, and the only stochasticity stems from the Gaussian mechanism.
\item \textbf{Very small learning rate.} If $\eta\ll\frac{1}{L}$, the optimization term dominates, yielding a slower $\mathcal{O}(1/(\eta T))$ decay, illustrating the necessity of balancing $\eta$.
\end{itemize}

\section*{Discussion}
The proof follows the standard smooth‑function descent analysis, with the only additional ingredient being Lemma \ref{lem:noise} that captures both sampling variance and the Gaussian mechanism. The clipping operation does not appear explicitly in the bound because it only guarantees that Assumption \ref{ass:clip} holds, which in turn ensures a finite variance $G^{2}$. The resulting rate matches the best known guarantees for DP‑SGD on non‑convex objectives and makes explicit how the privacy parameters $(\sigma,d)$ degrade utility.

\end{document}